\section{Scope and evidence procedure}
\label{sec:methods}
This structured review synthesizes a research collection covering literature available through 3 October 2026. The collection comprises source summaries, publication metadata, figure provenance, and verification notes linked to primary papers, developer documents, technical reports, and incident accounts. We reorganized and coded these materials rather than conducting a new exhaustive database search.

The collection's metadata table contains 153 records and its source summaries 147 detailed or short entries. These are inventory counts with different inclusion rules, not estimates of unique publications. A journal article and its preprint represent the same work unless the underlying study differs; the 8 October Nature Communications publication of the Piao study is therefore a status update, not an additional observation. Two targeted additions within the cutoff, the 30 September RHEON preprint and a 2 October paper on logit-derived contagion kernels, are discussed in Section~\ref{sec:collective}. They did not come from a new search.

We extract the unit of analysis, interaction mechanism, scale semantics, comparison condition, reported outcome, and qualifications from the existing source records. Numerical findings retain their stated denominator and budget condition. An estimate read from a source figure remains an estimate. A missing field is recorded as not reported in the extracted evidence, rather than inferred from a product name or an adjacent experiment. Appendix~\ref{app:evidence} applies the full schema to 57 selected evidence units. Appendix~\ref{app:inventory} separately lists all 153 metadata records with their source titles, evidence classes, and model/scale fields. The latter is an administrative inventory, not a second set of fully coded studies.

Evidence is separated into controlled agent experiments, benchmarks and datasets, developer system cards, engineering reports, observational studies, incident reconstructions, surrogate simulations, and theoretical or interpretive accounts. These classes describe what an observation can establish. A benchmark may isolate a causal contrast but have weak external validity; an incident may establish that a behavior occurred in deployment while offering no controlled estimate of its cause. Publication venue and evidence class therefore remain distinct.

The synthesis compares mechanisms within these boundaries. We do not combine task accuracy, speedup, dollar cost, agreement, or attack frequency into a common effect size. Neither a tally of positive results nor a cross-task average is interpreted as a universal law. The systematic component is the explicit coding framework and evidence mapping, not a claim of exhaustive retrieval or preregistered review methodology.
